\section*{Chapitre \ref{ch:b2:blood-pressure} --- La régulation de la pression artérielle et l'exercice}

\begin{solution}{exo:b2:blood-pressure:1}
$\bar P - P_{\text{ven}} = Q\,R = f\,V_s\,R$. La fréquence cardiaque (le
\omterm{prop:b2:heart:conduction}{nœud sinusal}, sous la commande du \omterm{prop:b2:heart:control}{nerf vague} et des nerfs
sympathiques)~; le volume d'éjection (le \omterm{def:b2:heart:anatomy}{ventricule}~: son remplissage,
fixé par les \omterm{def:b2:blood-circulation:vessels}{veines}, et sa contractilité, fixée par les nerfs
sympathiques)~; la résistance (le muscle lisse artériolaire, sous la
commande des nerfs sympathiques, des hormones et des métabolites
locaux).
\end{solution}

\begin{solution}{exo:b2:blood-pressure:2}
Capteurs~: récepteurs d'étirement des sinus carotidiens et de la crosse
aortique. Centre~: les centres cardiovasculaires du bulbe rachidien.
Effecteurs~: le \omterm{prop:b2:heart:control}{nerf vague} (fréquence) et les nerfs sympathiques
(fréquence, contractilité, tonus artériolaire et veineux). Rétroaction
négative, délai d'environ une seconde. Sans les capteurs, la pression
moyenne sur une journée reste normale, mais la pression varie
fortement d'une minute à l'autre à chaque changement de posture et à
chaque émotion.
\end{solution}

\begin{solution}{exo:b2:blood-pressure:3}
Le rein (cellules artériolaires) libère de la rénine lorsque sa
perfusion ou l'apport de sodium diminuent ou que ses nerfs sympathiques
s'activent~; la rénine découpe l'angiotensinogène (foie) en
angiotensine I~; une enzyme pulmonaire la convertit en angiotensine II,
qui contracte les \omterm{def:b2:blood-circulation:vessels}{artérioles} et stimule la corticosurrénale, laquelle
sécrète l'aldostérone~; l'aldostérone fait retenir par le rein du sodium
et de l'eau, ce qui élève le volume sanguin.
\end{solution}

\begin{solution}{exo:b2:blood-pressure:4}
La dilatation métabolique des \omterm{def:b2:blood-circulation:vessels}{artérioles} du muscle par ses propres
métabolites~; un \omterm{thm:b2:heart:starling}{débit cardiaque} multiplié par cinq sous l'effet des nerfs
sympathiques et de la pompe musculaire~; la constriction sympathique des
\omterm{def:b2:blood-circulation:vessels}{artérioles} de l'intestin, des reins et des muscles au repos, qui cèdent
du débit.
\end{solution}

\begin{solution}{exo:b2:blood-pressure:5}
$83\times 1.2 = \qty{100}{mmHg}$~; $333\times 0.35 = \qty{117}{mmHg}$.
Résistance divisée par $1.2/0.35 = 3.4$~: rayon multiplié par
$3.4^{1/4} =
1.36$.
\end{solution}

\begin{solution}{exo:b2:blood-pressure:6}
$30/(1 + G)$~: 10, 6 et \qty{3.3}{mmHg}. Avec un gain réduit de moitié,
la chute résiduelle double~: elle a des vertiges en se levant, voit gris
et peut s'évanouir.
\end{solution}

\begin{solution}{exo:b2:blood-pressure:7}
$\rho g h$~: $1060\times 9.81\times 0.4 = \qty{4160}{Pa} =
\qty{31}{mmHg}$~; $1060\times 9.81\times 1.2 = \qty{12500}{Pa} =
\qty{94}{mmHg}$. Cerveau $95 - 31 = \qty{64}{mmHg}$~; cheville
$95 + 94 =
\qty{189}{mmHg}$. La tête d'une girafe, à \qty{2.5}{m}, coûte
\qty{195}{mmHg}~: il lui faut une pression moyenne de quelque
\qty{250}{mmHg} au niveau du \omterm{def:b2:heart:anatomy}{cœur}, ainsi que des parois vasculaires
épaisses et une peau serrée dans les pattes.
\end{solution}

\begin{solution}{exo:b2:blood-pressure:8}
$Q = 280/(190 - 130) = \qty{4.7}{L/min}$. Athlète~: $5000/(200 - 30) =
\qty{29}{L/min}$~; volume d'éjection $\num{29400}/190 = \qty{155}{mL}$.
\end{solution}

\begin{solution}{exo:b2:blood-pressure:9}
Chute de $35/5 = \qty{7}{mmHg}$~: environ \qty{88}{mmHg}. Restauration
du volume~: réabsorption du liquide interstitiel dans les \omterm{def:b2:blood-circulation:vessels}{capillaires}
(de quelques minutes à une heure)~; réduction de l'urine sous l'effet de
l'\omterm{prop:b2:blood-pressure:raas}{hormone antidiurétique} et de l'aldostérone (de quelques heures à
quelques jours)~; soif et boisson (quelques heures)~; \omterm{def:b2:blood-circulation:blood}{hématies} produites
par la moelle osseuse (quelques semaines). La peau est pâle et froide
parce que le réflexe a fermé ses \omterm{def:b2:blood-circulation:vessels}{artérioles} pour réserver la pression au
cerveau et au \omterm{def:b2:heart:anatomy}{cœur}.
\end{solution}

\begin{solution}{exo:b2:blood-pressure:10}
Muscle à $R = 0.5$~: $95/0.5 = \qty{190}{mL/s}$, soit \qty{11.4}{L/min}~;
avec les \qty{4}{L/min} des autres organes, le \omterm{def:b2:heart:anatomy}{cœur} devrait fournir
\qty{15.4}{L/min}. Avec $Q$ fixé à \qty{5}{L/min}, la résistance totale
passe de 1.14 à $1/(1/1.14 - 1/5 + 1/0.5) =
\qty{0.37}{mmHg.s/mL}$ et la pression à \qty{31}{mmHg} --- le cerveau
défaillerait. L'organisme élève le débit vers le premier cas et contracte
les autres territoires, si bien que la pression se maintient et que le
muscle reçoit son débit.
\end{solution}

\begin{solution}{exo:b2:blood-pressure:11}
(a) Le rein mal irrigué a perçu une pression basse et libéré de la
rénine, et l'angiotensine et l'aldostérone ont élevé la résistance et le
volume. (b) Le \omterm{def:b2:blood-pressure:baroreflex}{baroréflexe} amortit les variations mais se réajuste sur
tout niveau qui persiste, et ne peut pas modifier la rétention de volume
par le rein. (c) La source de rénine, et l'organe qui réclamait une
pression plus élevée, avaient disparu. (d) Il aurait bloqué la
conversion en angiotensine II~: résistance plus faible, moins
d'aldostérone, pression plus basse.
\end{solution}

\begin{solution}{exo:b2:blood-pressure:12}
Le \omterm{def:b2:blood-pressure:baroreflex}{baroréflexe} agit en une seconde, avec un gain d'environ quatre et une
valeur de consigne qui dérive vers la pression dominante~: il amortit les
fluctuations et ne peut pas maintenir un niveau. Le rein agit sur des
heures et des jours, avec un gain en pratique infini --- il continue
d'excréter ou de retenir jusqu'à ce que la pression pour laquelle il est
réglé soit atteinte --- et sa consigne est fixée par sa propre
physiologie~: il maintient le niveau. Une élévation chronique signifie
donc que la consigne du rein s'est déplacée, et seuls des médicaments qui
agissent sur la boucle rénale ou sur la résistance qu'elle réclame la
font baisser.
\end{solution}

\begin{solution}{pb:b2:blood-pressure:1}
\textbf{1.} $1/R = 1/7.6 + 1/22.8 + 1/4.1 + 1/5.2 + 1/5.7 + 1/19 +
1/28.5 = 0.875$~: $R = \qty{1.14}{mmHg.s/mL}$~; $\bar P/Q = 95/83.3 =
1.14$.
\textbf{2.} $\bar P/R_i$~: cerveau 0.75, \omterm{def:b2:heart:anatomy}{cœur} 0.25, intestin et foie
1.39, reins 1.10, muscles 1.0, peau 0.30, autres \qty{0.20}{L/min}~;
fractions 15, 5, 28, 22, 20, 6 et \qty{4}{\%}.
\textbf{3.} $5000/70 = \qty{71}{mL}$~; $5\times 50 = \qty{250}{mL}$ de
dioxygène par minute.
\textbf{4.} Cerveau $0.75/1.4 = \qty{0.54}{L/min/kg}$~; reins $1.1/0.3
= \qty{3.7}{L/min/kg}$, sept fois celui du cerveau et cinquante fois la
moyenne du corps~: les reins sont irrigués pour filtrer, non pour se
nourrir.
\textbf{5.} Avec une pression unique partagée par tous les organes,
chacun prélève le débit que fixe sa propre résistance~; régler les débits
de façon centrale affamerait tout organe dont la demande changerait. La
pression est la monnaie commune.
\textbf{6.} \omterm{def:b2:heart:anatomy}{Cœur}~: $250\times 0.2\times 0.7 = \qty{35}{mL/min}$~;
muscle~: $1000\times 0.2\times 0.25 = \qty{50}{mL/min}$. Le \omterm{def:b2:heart:anatomy}{cœur} extrait
déjà la plus grande partie de ce qui passe~: il ne peut gagner du
dioxygène que par un surcroît de débit, et c'est pourquoi ses \omterm{def:b2:blood-circulation:vessels}{artérioles}
se dilatent dès qu'il travaille davantage.
\textbf{7.} $Q$ diminue de \qty{30}{\%} à $R$ fixe~: $\bar P$ diminue de
\qty{30}{\%}, soit \qty{28.5}{mmHg}, jusqu'à \qty{66.5}{mmHg}.
\textbf{8.} $28.5/5 = \qty{5.7}{mmHg}$~: environ \qty{89}{mmHg}.
\textbf{9.} $Q = 85\times 0.7\times 71 = \qty{4.2}{L/min} =
\qty{70}{mL/s}$~; $R = 89.3/70 = \qty{1.27}{mmHg.s/mL}$, soit
\qty{11}{\%} de plus.
\textbf{10.} $1060\times 9.81\times 0.4 = \qty{4160}{Pa} =
\qty{31}{mmHg}$~; cerveau $66.5 - 31 = \qty{35}{mmHg}$ avant le réflexe,
$89 - 31 = \qty{58}{mmHg}$ après.
\textbf{11.} $28.5/2 = \qty{14}{mmHg}$~: \qty{81}{mmHg} au niveau du
\omterm{def:b2:heart:anatomy}{cœur}, \qty{50}{mmHg} au niveau du cerveau --- vertiges, vision grise,
presque une syncope chaque fois qu'elle se lève d'une chaise.
\textbf{12.} La position allongée supprime la perte hydrostatique de
\qty{31}{mmHg} et ramène vers le \omterm{def:b2:heart:anatomy}{cœur} le \omterm{def:b2:blood-circulation:blood}{sang} accumulé~; le volume
d'éjection et la pression cérébrale se rétablissent en quelques
battements.
\textbf{13.} $45/5 = \qty{9}{mmHg}$~: \qty{86}{mmHg}.
\textbf{14.} $1/R = 0.875 - 1/19 - 1/4.1 + 1/38 + 1/8.2 = 0.727$~: $R =
\qty{1.38}{mmHg.s/mL}$. Cerveau $86/7.6 = \qty{11.3}{mL/s} =
\qty{0.68}{L/min}$~; peau $86/38 = \qty{2.3}{mL/s} = \qty{0.14}{L/min}$~;
intestin $86/8.2 = \qty{10.5}{mL/s} = \qty{0.63}{L/min}$.
\textbf{15.} Avec moins de \omterm{def:b2:blood-circulation:blood}{sang} et des \omterm{def:b2:blood-circulation:vessels}{artérioles} plus serrées, la
pression \omterm{def:b2:blood-circulation:vessels}{capillaire} passe sous la \omterm{thm:b2:blood-circulation:oncotic}{pression oncotique} sur toute la
longueur du \omterm{def:b2:blood-circulation:vessels}{capillaire}~: le liquide est réabsorbé au lieu d'être filtré~;
le \omterm{def:b2:blood-circulation:blood}{plasma} est dilué et l'\omterm{def:b2:blood-circulation:blood}{hématocrite} diminue.
\textbf{16.} L'aldostérone, par la rénine et l'angiotensine, déclenchée
par la faible perfusion du rein et l'activation sympathique~; l'\omterm{prop:b2:blood-pressure:raas}{hormone
antidiurétique}, déclenchée par la baisse du volume et l'élévation de la
concentration plasmatique.
\textbf{17.} $5\times 10^{12}$ cellules à $2.5\times 10^{6}$ par
seconde~: $2\times 10^{6}$ s, environ trois semaines.
\textbf{18.} En dessous d'environ \qty{60}{mmHg}, la courbe du réflexe
est plate~: les vaisseaux sont aussi serrés et le \omterm{def:b2:heart:anatomy}{cœur} aussi rapide qu'ils
peuvent l'être, le gain est nul, et chaque perte supplémentaire abaisse
la pression de toute sa valeur~; les tissus mal irrigués libèrent alors
de l'acide et se dilatent, et la pression s'effondre.
\textbf{19.} Muscles $110/0.33 = \qty{333}{mL/s} = \qty{20}{L/min}$~;
peau $110/5 = \qty{22}{mL/s} = \qty{1.3}{L/min}$.
\textbf{20.} Résistance du cerveau $8.8$ (débit 0.75 à 110), du \omterm{def:b2:heart:anatomy}{cœur}
$6.6$ (débit 1.0)~: $1/R = 1/8.8 + 1/6.6 + 1/8.2 + 1/10.4 + 1/0.33 + 1/5 +
1/28.5 = 3.75$~: $R = \qty{0.27}{mmHg.s/mL}$~; $Q = 110/0.27 =
\qty{410}{mL/s} = \qty{25}{L/min}$.
\textbf{21.} $\num{24700}/180 = \qty{137}{mL}$.
\textbf{22.} $24.7\times(200 - 40) = \qty{3950}{mL/min}$, seize fois les
\qty{250}{mL/min} du repos.
\textbf{23.} Débit $\times 4.9$, extraction $160/50 = \times 3.2$~;
$4.9\times 3.2 = 16$.
\textbf{24.} La commande du cortex moteur et les signaux des récepteurs
musculaires repondèrent, dans le bulbe, les informations des
\omterm{def:b2:blood-pressure:baroreflex}{barorécepteurs}, si bien que le réflexe défend \qty{110}{mmHg} au lieu de
95 --- la consigne est élevée, la boucle n'est pas désactivée.
\textbf{25.} Chute résiduelle en position debout \qty{5.7}{mmHg}~;
\qty{86}{mmHg} après l'hémorragie~; \qty{25}{L/min} et \qty{4}{L} de
dioxygène par minute à l'effort.
\end{solution}
